\section{Introduction}

Knowledge tracing (KT) estimates a learner's changing knowledge from a
chronological sequence of educational interactions
\citep{corbett-anderson-1995,abdelrahman-etal-2023}.  Although KT models vary in
their assumptions and architectures, their inputs normally associate each item
with one or more skills or \emph{knowledge components} (KCs).  A binary
Q-matrix makes this content model explicit: its $(i,k)$ entry records whether
item $i$ exercises KC $k$.  This matrix is often treated as supplied data, but
it already embodies a consequential hypothesis about the boundaries of
knowledge.

That hypothesis is especially difficult to specify for open-ended language
tutoring.  The Teacher--Student Chatroom Corpus version 2 (TSCCv2), for
example, contains 260 written English lessons, 41.4K conversational turns, and
approximately 363K word tokens from two teachers and thirteen learners.  Its
annotations include discourse sequences, response relations, pedagogical
focus, and grammatical-error corrections \citep{caines-etal-2022-teacher}.
These annotations make TSCCv2 a motivating future application for dialogue KT,
but they do not turn every learner utterance into a predefined assessment item.
A turn may contain explanation, scaffolding, elicitation, correction, or
several grammatical phenomena at once.  Neither the assessment opportunity nor
its KC decomposition is therefore given by the corpus.

Dialogue-based KT makes this missing representation visible.  \citet{scarlatos-etal-2025-dialogue}
use large language models (LLMs) to identify KCs and correctness at individual
turns in mathematics tutoring dialogues before applying KT models, with expert
annotations used to assess the automatic labels.  More recent work additionally
models the difficulty of tutor-posed tasks \citep{huang-etal-2026-interpretable}.
These studies operate in domains where an existing mathematical standard
provides a comparatively explicit candidate skill space.  Grammar raises an
earlier question: how should that space itself be constructed?

Consider the erroneous form \emph{Why didn't he went?}.  A coarse analysis
might assign one \emph{past-question} KC.  A factorised analysis might instead
identify finite-past selection, DO-support, negation placement, and
non-subject-WH formation.  An interaction-sensitive analysis might additionally
represent the dependency between DO-support and inversion.  These alternatives
share the same observed sentence but imply different learner states,
opportunity counts, and Q-matrices.  No decomposition is neutral.

Proficiency resources offer curriculum evidence but do not solve this problem
directly.  The Common European Framework of Reference for Languages (CEFR)
organises proficiency through broad functional scales \citep{coe-2020-cefr},
whereas the English Grammar Profile (EGP) associates more detailed grammatical
competence statements with CEFR levels \citep{okeeffe-mark-2017}.  An EGP
descriptor may nevertheless leave a grammatical property unspecified, cover
several alternatives, or instantiate the same structural configuration as a
different descriptor.  Treating each descriptor as a KC would conflate source
wording, curriculum evidence, grammatical structure, and a cognitive modelling
hypothesis.

\system separates these layers.  It first maps source descriptors into a
closed canonical representation, then constructs concrete grammatical items
without consulting a KC ontology.  Synthetic learner outcomes are generated
from a separate, declared structural oracle.  Only after both the item bank and
event stream have been fixed are candidate KCs selected or supplied, projected
onto items, and evaluated.  This order addresses an important confound: if a KC
policy determines which items exist or how responses are generated, predictive
differences between policies cannot be attributed to representation alone.

This report asks three questions:
\begin{enumerate}
  \item[RQ1] How reliably can EGP descriptors be normalised into a declared
        canonical representation of English verbal morphosyntax?
  \item[RQ2] Which candidate grammar distinctions are supported,
        identifiable, and retained using development-only structural evidence?
  \item[RQ3] How do frozen KC representations differ in Q-matrix coverage and
        transfer to unseen grammatical combinations under identical synthetic
        learner--item outcomes?
\end{enumerate}

The main contributions are: (i) a conservative six-dimensional
\grammarcell representation with explicit uncertainty and source provenance;
(ii) a fixed, split-independent controlled item bank; (iii) a development-only
KC selection procedure that diagnoses nuisance sensitivity and observational
equivalence; and (iv) an ontology-independent simulation and frozen-probe
protocol for controlled representation comparison.  These are methodological
and software contributions.  The study does not claim that the selected
dimensions are cognitively atomic, that the synthetic learners approximate
human acquisition, or that results transfer directly to conversational
language learning.

The remainder of the report reviews the relevant literature, formalises the
task and pilot data, describes the implemented pipeline and evaluation design,
states the present evidence boundary, and discusses the interpretation and
limitations of the framework.
